\documentclass[11pt,a4paper]{article}
\usepackage[margin=0.85in]{geometry}
\usepackage[T1]{fontenc}
\usepackage[utf8]{inputenc}
\usepackage{lmodern,microtype}
\usepackage{amsmath,amssymb,booktabs,tabularx,array,graphicx}
\usepackage{caption}
\usepackage[hidelinks]{hyperref}
\usepackage{enumitem}
\usepackage{float}
\usepackage{natbib}
\setlist{nosep,leftmargin=*}
\setlength{\parindent}{1.05em}
\setlength{\parskip}{0.10em}
\renewcommand{\arraystretch}{1.17}
\setlength{\emergencystretch}{2em}
\newcommand{\NA}{\textemdash}
\newcommand{\Pending}{\textit{Awaiting verified experiment}}
\newcommand{\Rtest}{D_{\mathrm{R}}^{\mathrm{test}}}
\newcommand{\Rtrain}{D_{\mathrm{R}}^{\mathrm{train}}}
\newcommand{\Rval}{D_{\mathrm{R}}^{\mathrm{val}}}
\newcommand{\Synthetic}{D_{\mathrm{S}}}
\title{\bfseries DermSynReal: Evaluating Morphology-Guided Synthetic Facial Images for Six-Class Dermatology Classification}
\author{Boda Shanmukha Datta\\
\small Department of Instrumentation and Control Engineering\\
\small Dr.\ B.\ R.\ Ambedkar National Institute of Technology Jalandhar\\
\small \textit{Additional contributors and clinical reviewers to be credited before submission}}
\date{\small Manuscript draft --- methodological sections complete; empirical analysis pending}
\begin{document}
\maketitle
\begin{abstract}
Clinical photographs of facial skin conditions vary widely in image quality, disease presentation and acquisition source. Classifiers trained on such collections can learn source-specific cues as readily as clinical morphology, making apparent improvements from synthetic augmentation difficult to interpret. This paper develops \emph{DermSynReal}, a controlled six-class study of whether morphology-guided diffusion images can improve classification of unseen real facial photographs. We define a diagnosis- and provenance-reviewed real-image collection, duplicate-aware group partitions, a structured prompt bank with isolated prompt families, and two matched generation strategies: generic disease-name descriptions and explicit lesion-morphology descriptions. We distinguish base diffusion synthesis from optional train-only low-rank adaptation and require independent review of generated-image morphology, artifacts and possible duplication. Five classifier pipelines---ResNet-50, EfficientNet-B3, ConvNeXt Tiny, DINOv2 ViT-S/14 and a PanDerm Base linear probe---are organized around fixed real validation and test sets. The primary outcome is six-class macro-F1 on the real test set, with paired group-bootstrap intervals, class-level metrics, source-stratified analyses and training-seed variability. We specify the contrasts that isolate the effects of augmentation, prompt detail and synthetic-image selection. The implementation and evaluation protocol are provided; numerical findings and clinical claims will be reported only after audited Kaggle execution. This design makes it possible to distinguish better-looking synthetic images from synthetic images that demonstrably help on real data.
\end{abstract}
\noindent\textbf{Keywords:} clinical facial photographs; dermatology; synthetic data; diffusion models; morphology-guided prompting; data leakage; transfer learning; reproducibility.
\par\medskip
\noindent\textbf{Status note.} This draft contains the complete methodological and contextual narrative. Tables in the Results section deliberately contain no invented measurements. References to planned experiments do not imply that a GPU run, label audit or clinical validation has already succeeded.
\section{Introduction}
A photograph of a facial skin condition contains more than a lesion. Its appearance depends on illumination, camera distance, skin texture, treatment history, facial location and the source from which the image was obtained. These differences matter to machine learning because a classifier has no inherent reason to attend to clinically meaningful features. When one disease category is collected mainly from one website or photographic setting, a model may recognize the source rather than the condition. An improvement on a random image split can therefore be misleading, particularly if resized versions of the same photograph appear in different partitions.

Clinical dermatology datasets face persistent problems of representation and source variation. The Fitzpatrick17k work drew attention to skin-type coverage and classification differences across skin types \citep{groh2021fitzpatrick}. SCIN adopted a different acquisition strategy, collecting user-contributed photographs of everyday skin concerns \citep{ward2024scin}. These efforts make valuable images available, but neither automatically supplies a diagnosis-verified, patient-independent six-class collection of facial photographs. A labeled image can depict a body site other than the face, more than one skin finding, or a diagnosis incompatible with the category to which it has been assigned.

Generative models offer a possible way to enrich the visual range of a limited dataset. Latent diffusion models synthesize images by denoising compressed representations under text conditioning \citep{rombach2022latent}. Yet a realistic generated face is not necessarily an informative dermatology sample. The prompt \emph{acne} leaves lesion count, size, inflammation and anatomical distribution unspecified. A more descriptive prompt can request those features, but it may also exceed the text encoder's effective conditioning capacity or produce incorrect lesions. Whether detailed prompts help a classifier is therefore a question for controlled evaluation, not an assumption about the generator's realism.

This paper presents \emph{DermSynReal}, an experimental framework designed around one central question: \emph{does morphology-guided synthetic-image augmentation change performance on unseen real facial photographs, beyond what is achieved with real-only learning or generic disease-name synthetic prompts?} We distinguish three effects. RQ1 concerns augmentation itself, comparing real-only training with real-plus-morphology training. RQ2 concerns prompt design, comparing generic and morphology-guided samples when the generator, sampling budget and learning schedule are matched. RQ3 concerns selection, testing whether quality- or diversity-based choice improves over a random matched-count sample from the same morphology pool. An optional synthetic-only arm tests the real-versus-synthetic domain gap.

The contribution is a reproducible \emph{evaluation design} accompanied by implementation rather than a claim of superior classification. It brings together provenance checks, duplicate-aware real-image grouping, a structured prompt library, base diffusion and optional training-only LoRA adaptation, five visual encoders, and paired evaluation on a frozen real test. This separation matters: a benefit attributed to clinical prompt detail is confounded if the generic and detailed images come from different generator checkpoints. Similarly, a selection gain cannot be isolated if one arm contains more training examples or was tuned on the real test.

The scope is a six-output photographic research classifier, not medical diagnosis, treatment guidance or prospective clinical validation. The intended labels contain clinically important subtypes, especially within fungal, pigmentary and malignant or premalignant conditions. Final conclusions must acknowledge uncertain diagnoses, image licenses, unknown patient identities and the limits of a single collected test domain. Numerical conclusions are left open until those questions have been addressed through actual runs and review.

\section{Related Work}
\subsection{Clinical datasets, representation and shortcuts}
Dermatology collections have enabled transfer learning but have also exposed variability in annotation and acquisition source. Fitzpatrick17k uses clinical photographs annotated for skin type, highlighting the importance of representation in evaluating learning systems \citep{groh2021fitzpatrick}. SCIN broadens the setting through user-contributed images and condition information \citep{ward2024scin,scin_repository}. Neither removes the need to check anatomical site, original diagnosis, photographic modality and permitted use when assembling a new task. Different crops of a single photograph, watermarks or site-specific presentation can also create shortcuts. If the same visual subject reaches both training and testing, a classifier may appear to generalize while recognizing an already-seen image. We therefore consider label and leakage auditing part of the research design, not merely data cleaning.

\subsection{Diffusion-generated dermatology images}
Latent diffusion offers a flexible approach to text-conditioned synthesis \citep{rombach2022latent}. Dermatology research has begun exploring whether generated lesions help train classifiers; Derm-T2IM studies synthetic \emph{dermoscopic} lesions with vision transformers and CNNs \citep{farooq2024dermt2im}. A dermoscopic crop differs from full-face photography in its anatomical context, image scale and acquisition artifacts. Results from one modality are not evidence that generated faces improve another. Moreover, a model may synthesize a plausible skin texture while omitting the requested lesion, distorting its border or introducing a consistent synthetic artifact. In this study, lesion fidelity, image novelty and classification utility are deliberately measured separately. An aesthetic score or embedding similarity alone is not treated as a clinical truth label.

\subsection{General and domain-specialized encoders}
ResNet-50 supplies residual convolutional learning \citep{he2016resnet}; EfficientNet-B3 represents compound scaling \citep{tan2019efficientnet}; and ConvNeXt Tiny modernizes convolutional blocks while remaining a CNN \citep{liu2022convnext}. DINOv2 ViT-S/14 provides features learned by large-scale self-supervision \citep{oquab2024dinov2}. These families offer distinct feature representations for testing whether an augmentation effect transfers across classifiers rather than appearing on one backbone only.

PanDerm is especially relevant as a dermatology-pretrained foundation model trained across multiple imaging modalities \citep{yan2025panderm}. Its published results motivate examination of specialty representations but do not establish superiority on our six-category facial task. Our PanDerm Base experiment is a frozen-feature linear probe; it must not be described as full fine-tuning. The four general-purpose backbones instead support head training, final-stage unfreezing and optional full fine-tuning.

\subsection{Unresolved methodological question}
The gap addressed here is not the absence of synthetic dermatology work. It is the difficulty of attributing an improvement to a specific intervention if prompt detail, generator weights, synthetic-image selection, classifier optimization and test composition all change simultaneously. DermSynReal separates the augmentation, prompt and selection effects into explicit comparisons. When an intended match is not feasible, the result is described as a confounded exploratory comparison rather than evidence that morphology wording caused the difference.

\section{Task Definition and Real-Image Data}
\subsection{Output taxonomy}
We study a single-label, six-category classification task for clinical-looking facial photographs. The proposed output vocabulary is
\begin{equation}
\mathcal{Y}=\{\mathrm{Acne},\mathrm{Fungal},\mathrm{Vitiligo},
\mathrm{Hyperpigmentation},\mathrm{SkinCancer},\mathrm{Normal}\}.
\label{eq:labels}
\end{equation}
These categories describe the research endpoint, not a comprehensive dermatological differential diagnosis. Clinically important subtypes remain separate in the source metadata and generation specification. Acne may involve comedonal or inflammatory appearances; hyperpigmentation encompasses distinct etiologies; facial fungal conditions require diagnosis-specific review. The skin-cancer category needs particularly careful definition: melanoma, basal-cell carcinoma and squamous-cell carcinoma are malignancies, whereas actinic keratosis is premalignant. Whether actinic keratosis belongs in this broad output must be justified and fixed before testing. A photograph with uncertain or multiple diagnoses is not simply forced into the nearest class.

An eligible real-image record is written as $r_i=(x_i,y_i,s_i,\ell_i,g_i,h_i,a_i)$, comprising an image $x_i$, mapped label $y_i$, collection source $s_i$, original diagnostic provenance $\ell_i$, protective group $g_i$, SHA-256 content fingerprint $h_i$ and recorded approval decision $a_i$. The original source descriptor and its mapped class must be compatible; a folder name alone is not proof.

\subsection{Acquisition and review}
The implemented pipeline targets one Kaggle real-image collection, version 2, under the dataset slug \texttt{shanmukhadattaboda/facial-skin-conditions-real-image-collection}. Its six raw folders are discovered with their original case-sensitive names. Uploaded image counts represent candidate files rather than verified unique patients or medically established diagnoses. Original publication and reuse rights must be checked separately: a license selected by a dataset uploader cannot automatically override copyright or restrictions on photographs from another source.

For every candidate we retain a stable ID, original relative path, source, original diagnosis when available, reviewer ID, decision, reason for rejection and content digest. The workflow flags obvious contradictions for manual inspection, such as Stevens--Johnson syndrome or toxic epidermal necrolysis photographs placed in a skin-cancer folder, or rosacea and dermatitis images placed under normal skin. These flags are screening aids, not an automated dermatologist. Images lacking credible diagnostic provenance, adequate rights or the intended facial presentation are rejected or held for review. Accepted labels are then frozen before partitioning.

\subsection{Exact and near-duplicate grouping}
Exact duplication is screened by both original file SHA-256 and decoded-pixel hashes. Potential near-duplicates are proposed using flip-aware perceptual hashes, additional gradient hashes and pretrained visual features; ORB-based geometric matching and a gallery aid ambiguous review. For pHash, one candidate distance is
\begin{equation}
d_{\rm ph}(i,j)=\min_{u,v\in \mathcal{T}}
d_H\!\left({\rm pHash}(u(x_i)),{\rm pHash}(v(x_j))\right),
\qquad
\mathcal{T}=\{{\rm id},H,V,R_{180}\},
\label{eq:phash}
\end{equation}
where $d_H$ is Hamming distance. A threshold of approximately five bits is a candidate heuristic rather than a diagnostic assertion of image identity. Candidate links are consolidated with connected components, and conflicting-label matches are quarantined. The transitivity of these groups can link ambiguous photographs; such cases deserve review rather than automatic reinterpretation as the same patient.

The protective constraint is
\begin{equation}
g_i=g_j\quad\Longrightarrow\quad
{\rm split}(x_i)={\rm split}(x_j).
\label{eq:group}
\end{equation}
Actual patient or encounter IDs, if legitimately available, take precedence over image-only group inference. Otherwise the manuscript will describe the resulting partitions as \emph{visually grouped}, not as proven patient-disjoint.

\subsection{The immutable real split}
The planned real partition proportions are approximately 70\% training, 15\% validation and 15\% test. Indivisible groups and imbalanced categories may prevent exact fractions; actual counts will be reported. Let $\Rtrain$, $\Rval$ and $\Rtest$ denote these frozen datasets. The test set is excluded from image selection, generator LoRA adaptation, prompt revision, threshold tuning, classifier training and checkpoint selection. The validation set selects classifiers and training settings but is never used to fit the generator.

The same globally grouped split is reused across the six LoRA inputs. A class-specific LoRA receives only its label's \emph{training} subset; it may report counts for that class in validation and test without fitting on those photographs. Re-splitting each raw disease folder independently could break group protection across labels and is deliberately prohibited.

\subsection{Limits of the dataset safeguards}
These checks make known label conflicts and duplicate risks visible, but they do not guarantee diagnostic truth, complete near-duplicate coverage or patient independence. They also do not establish that the source is representative of the clinical population. Dataset inclusion totals, exclusions, class support, duplicate groups and reviewer decisions are therefore reserved for the measured Results section rather than assumed from the uploaded image listing.

\section{Morphology-Guided Synthetic Image Generation}
\subsection{Structured, family-aware prompts}
The six generator directories are linked to an authored prompt collection. Each disease category contains 300 prompt records grouped into 60 prompt families. Within each category, 210 prompts belong to the training partition, 45 to validation and 45 to test; all instances of a family remain in one partition. This allocation makes it possible to examine generalization to unseen prompt formulations without confusing that diagnostic with the primary real-image test.

A morphology record $p$ specifies a set of visual attributes $\mathbf{m}$: visible lesion or pigment characteristics, spatial distribution on the face, approximate severity, local texture, skin appearance, camera viewpoint and illumination. For example, an acne prompt may distinguish scattered inflammatory papules on the cheek from a small number of noninflamed comedones on the forehead. A vitiligo prompt may request depigmented patches of differing extent and border, while a pigmentation prompt may describe bilateral malar patches rather than simply asking for a dark mark. These are instructions to an image generator, not statements that the generated findings are clinically verified. The hyperpigmentation, fungal and skin-cancer prompt families similarly preserve subtype information without implying that all subtypes are interchangeable.

For the prompt-design contrast, each morphology prompt has a corresponding generic counterpart. The generic version retains basic photographic context but names the broad condition without the detailed lesion attribute bundle. Both versions preserve the prompt ID, prompt-family allocation and, where feasible, the generator revision and seed schedule. We formalize the intended comparison as
\begin{equation}
x_{i,v,k}=G_{\theta,\phi}(p_{i,v},z_{i,k};\gamma),
\qquad v\in\{\mathrm{generic},\mathrm{morphology}\},
\label{eq:generator}
\end{equation}
where $G_{\theta,\phi}$ denotes the generator with base parameters $\theta$ and optional adaptation $\phi$, $z_{i,k}$ is the stochastic seed/noise choice and $\gamma$ collects guidance, sampling steps, resolution and scheduler settings. A difference between prompt variants can be associated with the wording only when the remaining inputs are adequately matched.

\subsection{Diffusion architecture and checkpoint identity}
The present generation scripts use an SD2.1-compatible base for the acne, fungal, vitiligo, hyperpigmentation and skin-cancer appearance tasks, and an SDXL-compatible pipeline for the normal-skin category. Both belong to the wider latent diffusion family \citep{rombach2022latent}, but they use different architectures and encoding procedures. SDXL output must not be presented as an isolated prompt-text comparison against SD2.1 output from another class. For each run we record the precise base repository, model revision or checkpoint hash, scheduler, sampler settings, positive and negative text, guidance, image size and seed.

Some external disease resources are textual-inversion concepts rather than complete diffusion models or LoRA adapters. We distinguish these mechanisms in the metadata. A condition represented by a dermatology concept learned at another anatomical site does not constitute a validated facial generator. Base-only pilots are therefore the appropriate first test before optional embeddings or adapters are enabled.

\subsection{Token limits and long-prompt conditioning}
A long prompt does not necessarily provide more effective control. The relevant text encoders have finite context windows, and silently truncating the final clauses may remove the very morphology being studied. Generator notebooks consequently validate prompt length and, where supported, experiment with Compel-based conditioning over bounded chunks. We do not assume that concatenating or chunking text guarantees faithful rendition of all attributes. The audit retains the effective token length, conditioning settings and rejected prompts so that apparent differences between generic and morphology arms are not caused by undocumented truncation.

\subsection{Optional LoRA adaptation}
Low-rank adaptation modifies selected weight matrices through trainable low-rank terms while leaving the original model weights fixed \citep{hu2022lora}. In simplified form,
\begin{equation}
W'=W+\Delta W,\qquad \Delta W=BA,\qquad
A\in\mathbb{R}^{r\times d},\ B\in\mathbb{R}^{d'\times r},
\end{equation}
with $r$ much smaller than the original dimensions. Its value here is practical: a disease-specific generator can be adapted without full diffusion-model retraining. The exact targeted modules, rank, optimization steps and checkpoint format must nevertheless be preserved because LoRA variants are not equivalent to textual inversion or to the unadapted base.

Only the reviewed real \emph{training} images are eligible for LoRA input. The six LoRA notebooks are wired to class-specific subsets derived from a single global frozen split, and a requested subtype requires a documented per-image subtype decision. Optional LoRA-generated examples form a separate, versioned generator source; they may be combined into a reviewed preprocessing pool but not silently mixed with base-model examples in a matched prompt experiment. A generator trained on real images creates additional memorization risk, so generated samples require novelty checks against the original real training set and especially against protected validation and test sets.

\subsection{Candidate counts, screening and provenance}
The generation design can produce five samples per authored prompt. Across the full 300-prompt bank that corresponds to 1,500 \emph{candidate} images per class: 1,050 from training families and 225 each from validation and test families. This is an intended production count before failures or review, not a claim of 1,500 accepted images. Only accepted images from training prompt families may be offered for classifier augmentation. The held-out synthetic prompt partitions are reserved for optional separate diagnostic analyses.

Every generated image receives its original prompt identifier, family, seed, generator/adapter identity, original image digest, status and review metadata. Human review considers whether the requested visible attributes are present, whether the face is plausible, whether watermarks or text are present, whether a second face or major artifact intrudes, and whether there are exact or near-copy concerns. An initially produced image has status \emph{needs review}; model confidence or image realism does not approve it. The resulting accepted pool is subjected to exact SHA and decoded-pixel hashing and to prompt-family leakage checks before being used in a classifier manifest.

\begin{figure}[ht]
\centering
\setlength{\fboxsep}{7pt}
\begin{tabular}{c}
\fbox{\parbox{0.83\linewidth}{\centering \textbf{Reviewed real facial photographs}\\
source/diagnosis checks $\rightarrow$ duplicate groups $\rightarrow$ fixed train/val/test}}\\[4pt]
$\downarrow$\\[-2pt]
\fbox{\parbox{0.83\linewidth}{\centering \textbf{Training-only generation}\\
authored generic and morphology prompt families; base diffusion; optional train-only LoRA}}\\[4pt]
$\downarrow$\\[-2pt]
\fbox{\parbox{0.83\linewidth}{\centering \textbf{Synthetic quality/novelty gate}\\
manual appearance assessment; duplicate quarantine; prompt-family preservation}}\\[4pt]
$\downarrow$\\[-2pt]
\fbox{\parbox{0.83\linewidth}{\centering \textbf{Controlled classifiers and evaluation}\\
E0--E5 training; real validation selection; identical frozen real test; paired statistics}}
\end{tabular}
\caption{Study workflow, showing the separation between generator construction and the held-out real-image evaluation. This is a methodological diagram, not an empirical result.}
\label{fig:workflow}
\end{figure}

\section{Classifier Training and Experimental Conditions}
\subsection{Five classification pipelines}
We use a fixed six-output prediction head and compare five available pretrained representations. The four general visual architectures are ResNet-50, EfficientNet-B3, ConvNeXt Tiny and DINOv2 ViT-S/14. A fifth pipeline uses genuine PanDerm Base pretrained features with a linear classifier. In particular, EfficientNet-B3 is not EfficientNet-B0, ConvNeXt Tiny is not ConvNeXt V2 Tiny, and DINOv2 ViT-S/14 is not the larger ViT-B/14. The manuscript, tables and checkpoint names must use only the models actually executed.

For a classifier $f_\omega(x)\in\mathbb{R}^{6}$, the objective is standard multiclass cross-entropy over the selected \emph{training} rows:
\begin{equation}
\mathcal{L}_{\rm CE}(\omega)=
-\frac{1}{N}\sum_{i=1}^{N}
\log\!\left(
\frac{\exp f_\omega(x_i)_{y_i}}
{\sum_{c=1}^{6}\exp f_\omega(x_i)_c}
\right).
\label{eq:loss}
\end{equation}
The output label order is fixed across training, checkpoint reload and evaluation. Resizing preserves photographic aspect ratio through a documented square-padding policy; images are decoded and converted to RGB, and the model's expected normalization is used. On-the-fly random augmentation belongs to training only. Neither validation nor test images receive stochastic augmentation, and corrupt or near-blank inputs are reported rather than silently assigned a class.

The general backbones follow three possible stages. First, the pretrained feature extractor is held fixed while the six-class head is learned. Second, final feature blocks are unfrozen at a smaller learning rate. Third, full backbone fine-tuning is optional; its use depends on recorded GPU resources and real-validation behavior. Optimizer parameter groups separate head and backbone rates. Early stopping and checkpoint selection use macro-F1 on $\Rval$, never $\Rtest$. A successful experiment must demonstrate that the saved checkpoint can be reloaded to reproduce the stored class ordering and predictions. The PanDerm pipeline is different by design: the official pretrained encoder remains frozen and a linear probe is fitted to extracted features. Consequently, its training cost and adaptation capacity cannot be assumed comparable to end-to-end fine-tuning.

\subsection{Controlled augmentation arms}
Table~\ref{tab:arms} defines the arms used to isolate the interventions. In every arm except optional synthetic-only E5, the same reviewed real training subset is retained; every arm is selected using the same real validation and tested on the same real test.

\begin{table}[ht]
\centering\small
\caption{Prespecified experiment arms. E4 and E5 are optional; E3 and E4 require measured selection metadata.}
\label{tab:arms}
\begin{tabularx}{\linewidth}{@{}lX X@{}}
\toprule
Arm & Training samples & Comparison addressed\\\midrule
E0 & Reviewed real training only & Real-only reference\\
E1 & Real + generic-name synthetic samples & Generic augmentation effect\\
E2 & Real + morphology-prompt synthetic samples & Morphology and augmentation effects\\
E3 & Real + quality-ranked samples from the E2 morphology pool & Selection by documented quality\\
E4 & Real + diversity-selected samples from the same morphology pool & Selection by documented coverage\\
E5 & Synthetic morphology samples only & Optional real-domain transfer diagnostic\\
\bottomrule
\end{tabularx}
\end{table}

The planned class budget is $n_s$ accepted synthetic training images per category, held equal in E1 and E2, and likewise in the E2/E3/E4 comparison. If a class lacks adequate accepted examples, the matched experiment is deferred or the budget is reduced for \emph{all} arms; it is not supplemented from held-out prompt families. For quality selection, a documented morphology score must be provided for candidate images. Diversity selection requires a cluster or coverage identifier. Neither score is derived from real-test performance. We record the selected synthetic IDs and prompt-family IDs for every generated manifest.

The E1--E2 contrast is valid as a prompt intervention only for matching generator architecture and adaptation mode, including checkpoint and adapter revision. The preprocessing pipeline may combine reviewed base and LoRA images into one stored pool, but the experiment-manifest builder filters by generator mode to make controlled arms. Without this filter, a gain cannot be attributed to prompt detail. Likewise, the different SDXL base used for normal-skin generation is a potential class-source shortcut: comparisons within each class should keep the generator identity stable, and the discussion must acknowledge that synthetic style could still correlate with labels.

\subsection{Optimization and reproducible comparison}
Each configured backbone uses an identified initialization, image size, batch size, optimizer, training-stage schedule and seed. The current manifest protocol reserves the training seeds 13, 42 and 2026 for repeated selected experiments. For timm classifiers, the implementation aims to match the number of optimizer updates per epoch to the real-only baseline even when an augmented pool is larger. This controls one component of training exposure, but it does not guarantee identical numbers of distinct real and synthetic examples seen or equal effective class balance. Such differences must be recorded, not hidden.

Where possible, the same initial checkpoint and random-seed convention are used for the corresponding E0/E1/E2 training runs. The real validation set selects the early-stopping checkpoint independently for each condition. Hyperparameter changes that occur after seeing real-test metrics constitute exploratory reanalysis; they cannot be presented as prespecified confirmatory results.

\subsection{Synthetic-only and diagnostic evaluation}
The optional E5 arm fits a classifier on accepted synthetic images from training prompt families while retaining real validation and real test for comparison. The evaluation of E5 on real images measures transfer from synthetic training to the collected real-image domain. If an independent synthetic prompt-test set was also generated, its performance is reported in a separate table. A combined real-plus-synthetic score must not pool real-test and synthetic-test predictions because the two domains have different meanings and independence assumptions.

An optional low-data study could restrict the real classifier training set to a subset of $\Rtrain$ while leaving $\Rval$ and $\Rtest$ untouched. If a disease-specific LoRA was adapted using the \emph{full} real training set, such a comparison would test limited classifier labels, not a truly low-real-data generation system. A strict low-data claim requires that all dataset-specific generator adaptation obey the smaller real-data budget.

\section{Evaluation and Statistical Analysis}
\subsection{Primary endpoint and class-level measures}
The primary endpoint is macro-averaged F1 on $\Rtest$. For each of the six categories $c$, precision $P_c$ and recall $R_c$ are calculated from the real-test confusion matrix. With an undefined class F1 taken as zero,
\begin{equation}
F_{1,\mathrm{macro}} =
\frac{1}{6}\sum_{c=1}^{6}
\frac{2P_cR_c}{P_c+R_c}.
\label{eq:f1}
\end{equation}
This definition gives equal weight to all six target categories regardless of their frequencies. We also report overall accuracy, balanced accuracy (mean class recall), each class's support and confusion matrix. A macro average is uninformative if one of the intended classes is absent from the test set; all six real-test class supports must therefore be recorded.

All classifiers must output scores in the same frozen class order. For one run, the prediction CSV preserves the test image ID, true and predicted labels, real image source, protective group, and six class probabilities when available. The test partition is evaluated only after the corresponding training configuration has been finalized using real validation.

\subsection{Paired real-test contrasts}
The main contrasts are
\begin{align}
\Delta_{\mathrm{aug}}&=F_1(E2)-F_1(E0),\\
\Delta_{\mathrm{prompt}}&=F_1(E2)-F_1(E1),\\
\Delta_{\mathrm{quality}}&=F_1(E3)-F_1(E2),\\
\Delta_{\mathrm{diversity}}&=F_1(E4)-F_1(E2).
\label{eq:contrasts}
\end{align}
All differences refer to \emph{the same real-test IDs}, within one backbone and training seed. We do not compare F1 differences computed on independently redrawn test sets or silently pair mismatched model checkpoints. E3 and E4 are additional selection contrasts, included only when their prerequisites are satisfied. E5 is diagnostic rather than an additional primary superiority claim.

To characterize test-sample uncertainty, the evaluation code forms paired bootstrap replicates of each F1 difference, resampling protective real-image groups where that is the available unit. If the groups were visually inferred rather than linked to verified patients, this limitation is disclosed. A percentile interval is reported as
\begin{equation}
{\rm CI}_{95\%}(\Delta)=
\left[Q_{0.025}\bigl(\Delta^{(b)}\bigr),
Q_{0.975}\bigl(\Delta^{(b)}\bigr)\right],
\label{eq:bootstrap}
\end{equation}
where each replicate evaluates both trained models on the same resampled test groups \citep{efron1986bootstrap}. The interval describes uncertainty under that resampling scheme; it is not proof of causality or an automatically valid p-value.

A separate summary of training seeds reports the mean, sample standard deviation (if at least two seeds are available), median and range of observed real-test scores. A new seed changes training randomness but does \emph{not} create an independent clinical test cohort. Accordingly, seed variability is shown separately from paired test-group bootstrap intervals. The current implementation does not manufacture p-values. If inferential hypothesis testing is subsequently preregistered, the primary family and appropriate multiplicity control, such as Holm's procedure \citep{holm1979multiple}, must be specified before interpreting significance.

\subsection{Source sensitivity, calibration and errors}
A classifier can perform differently across original collections even when its pooled macro-F1 appears stable. We therefore compute source-stratified counts and metrics where class support is sufficient, while marking a source with too few categories as a partial diagnostic rather than a complete six-class comparison. Per-class errors are grouped for review, with no patient-identifiable example published without permission.

Where six valid probability outputs exist, secondary analysis includes multiclass negative log-likelihood, Brier score and expected calibration error. Calibration is calculated by evaluation domain rather than pooling real and synthetic samples. Any temperature or threshold parameter is fitted using real validation rather than the real test. Class-specific and source-specific findings are descriptive when support is small.

\subsection{Interpretation rules}
The sign of one observed metric difference does not determine the study's conclusion. If E2 exceeds E0 but not E1, there may be an augmentation benefit without a morphology-specific benefit. If synthetic images receive high visual quality ratings but real-test performance does not improve, the generator has not shown downstream utility. If an apparent gain occurs mainly in one data source, source association may be a competing explanation. All results are reported with actual accepted-image counts, label-audit state, train budgets and checkpoint identifiers. Results from an unreviewed dataset cannot be promoted to clinical evidence simply by adding a confidence interval.

\section{Results}
\label{sec:results}
\noindent\textbf{This section is intentionally unpopulated.} Tables and figure slots below are the required structure for measured Kaggle outputs. Dashes indicate \emph{not yet measured}, not a zero score or a nonsignificant finding. The final text must be written after the dataset label audit and evaluation registry pass their checks. No quantitative outcomes are inferred from the software implementation.

\subsection{Real-data eligibility and frozen partitions}
\begin{table}[ht]
\centering\small
\caption{Final diagnosis-reviewed real-image inventory. Replace every dash with recorded reviewed counts and state the number of protective groups.}
\label{tab:realdata}
\begin{tabular}{@{}lrrrr@{}}
\toprule
Target class & Approved & Train & Val & Real test\\\midrule
Acne & \NA & \NA & \NA & \NA\\
Fungal infection & \NA & \NA & \NA & \NA\\
Vitiligo & \NA & \NA & \NA & \NA\\
Hyperpigmentation & \NA & \NA & \NA & \NA\\
Skin cancer/precancer scope & \NA & \NA & \NA & \NA\\
Normal & \NA & \NA & \NA & \NA\\\midrule
Total & \NA & \NA & \NA & \NA\\
\bottomrule
\end{tabular}
\end{table}
\noindent\Pending. Insert exclusion counts, exact/near-duplicate groups, cross-class conflicts, group-origin information and source-rights decisions. Explain if known patient IDs were unavailable.

\subsection{Generation yield and morphology review}
\begin{table}[ht]
\centering\small
\caption{Generator-prompt yield, accepted appearance quality and rejected duplicates. Each row must identify the actual generator revision and adaptation mode.}
\label{tab:synthetic}
\begin{tabularx}{\linewidth}{@{}Xrrrr@{}}
\toprule
Configuration & Candidates & Accepted & Accept (\%) & Near-copy rejects\\\midrule
E1 generic prompts & \NA & \NA & \NA & \NA\\
E2 morphology prompts & \NA & \NA & \NA & \NA\\
Optional LoRA source (separately) & \NA & \NA & \NA & \NA\\
\bottomrule
\end{tabularx}
\end{table}
\noindent\Pending. Insert a permitted representative-image grid only after originality, privacy, attribution and morphology review. The figure must distinguish real and synthetic images unambiguously.

\subsection{Primary real-test classification}
\begin{table}[ht]
\centering\small
\caption{Primary real-test performance by experiment arm. Make one version per backbone or add a backbone column when measured. Report mean across seeds and its SD separately.}
\label{tab:primary}
\begin{tabular}{@{}lrrr@{}}
\toprule
Arm & Macro-F1 & Balanced accuracy & Accuracy\\\midrule
E0 real only & \NA & \NA & \NA\\
E1 generic synthetic & \NA & \NA & \NA\\
E2 morphology synthetic & \NA & \NA & \NA\\
E3 quality selected & \NA & \NA & \NA\\
E4 diversity selected (optional) & \NA & \NA & \NA\\
E5 synthetic only (optional) & \NA & \NA & \NA\\
\bottomrule
\end{tabular}
\end{table}
\noindent\Pending. State the tested backbone, exact frozen real-test sample size, number of distinct protective groups, matched synthetic budget, training seeds, checkpoint version and failure/early-stopping events.

\subsection{Matched real-test contrasts}
\begin{table}[ht]
\centering\small
\caption{Paired differences in real-test macro-F1, evaluated on the identical real test for each backbone and seed. A CI is not a p-value.}
\label{tab:paired}
\begin{tabularx}{\linewidth}{@{}lXr@{}}
\toprule
Contrast & Interpretation & Difference [95\% group-bootstrap CI]\\\midrule
E2 versus E0 & Augmentation versus real only & \NA\\
E2 versus E1 & Morphological specificity versus generic prompts & \NA\\
E3 versus E2 & Quality selection at matched sample count & \NA\\
E4 versus E2 & Diversity selection at matched sample count & \NA\\
\bottomrule
\end{tabularx}
\end{table}
\noindent\Pending. Insert paired bootstrap intervals, observed signs across training seeds and a statement on the reliability of the underlying grouping unit. Do not add a fabricated significance column.

\subsection{Class errors, source behavior and calibration}
\noindent\Pending. Insert real-test six-class confusion matrices, per-class precision/recall/F1, source-specific support and calibrated probability diagnostics. Figures should come from the saved Evaluation 04--07 artifacts, with sources shown only at justified support thresholds. If synthetic prompt-test results exist, place them in a separate supplementary panel rather than pooling them into the primary real-test score.

\subsection{Result-dependent synthesis}
\noindent\Pending. Replace this single paragraph with a factual summary answering RQ1--RQ3, naming the backbones, effect magnitudes, uncertainty intervals, disagreement across seeds, and negative results. A clinically favorable interpretation must be justified by the data and cannot rest on visual quality alone.

\section{Discussion}
\subsection{What the comparisons can establish}
The experimental structure was chosen to make a narrow but useful conclusion possible. If the morphology arm performs better than real-only training, the result supports an augmentation effect for the reviewed data, the chosen generator and the tested classifier configuration. It would not establish that explicit morphology caused the improvement unless the morphology arm also differs from the matched generic arm. Conversely, if generic and morphology images produce similar scores, increasing prompt detail may have no incremental value even when both improve over E0. These distinctions are important because generation quality, dataset diversity and model optimization are related but different quantities.

Quality and diversity selection answer another question. If E3 differs from E2 when both draw the same number of images from the same morphology pool, a plausible explanation is the selected samples rather than prompt formulation. That explanation still needs caution: a reviewer may favor familiar or visually conspicuous lesions, making a quality-ranked subset less representative of ordinary clinical presentations. Diversity-based selection can also overemphasize unusual but low-fidelity candidates. The selection policy must therefore be reported with the acceptance criteria and classwise sample counts, including cases where it fails to improve real-test performance.

\subsection{The synthetic-to-real gap}
A synthetic-only classifier may learn visually consistent lesion patterns yet perform poorly on real photographs. This outcome is not contradictory. Generated images can have cleaner backgrounds, similar face framing, recurring skin texture or category-specific color artifacts absent from the target collection. Combining real and synthetic data may reduce such mismatch, but it may also cause the network to rely on generator-specific shortcuts. A source-separated confusion analysis is therefore more informative than declaring a generator successful because its images look medically plausible.

A gain that appears only on photographs from one source is a warning sign. The study can describe that effect, but it should not call it evidence of broad clinical generalization. The collected real test is not equivalent to an independent prospective hospital cohort. A stronger external-validity claim would require a separately sourced, diagnosis-reviewed set acquired under a different imaging process and frozen before tuning.

\subsection{Threats to internal validity}
Several practical limitations are anticipated. First, the broad six-category taxonomy simplifies clinical distinctions. A premalignant lesion and an invasive carcinoma should not be conflated without explicit endpoint design. Second, labels imported from filenames or public galleries can be wrong or incomplete. The review gate reduces known errors but cannot replace qualified diagnostic assessment. Third, source and watermark artifacts may remain correlated with labels. Fourth, visual near-duplicate grouping cannot fully establish independence among patients, and incomplete grouping may leave correlated images across the split.

Fifth, base generation and LoRA adaptation are different interventions. If their outputs are mixed, any classifier difference may reflect the adapter, the text prompt or the sample-selection policy. The analysis must preserve the separate generator identity. Sixth, normal-skin synthesis uses SDXL while most disease generations use SD2.1; this architecture difference may introduce a category-correlated synthetic style that resembles a class shortcut. Finally, training-seed variation and uncertainty from a limited real test are different sources of variability. Repeating training three times does not create three independent test sets.

\subsection{Scope of clinical interpretation}
Even an improvement supported by paired real-test intervals would not justify using the classifier to diagnose an individual. Clinical use requires reliable diagnoses, clinically meaningful differential endpoints, diverse patient populations, error analysis with qualified experts and prospective testing under realistic acquisition conditions. Subgroup statements require dependable annotations and sufficient representation; appearance cannot be used to infer sensitive personal traits. The study is primarily a controlled analysis of synthetic augmentation under a specified photographic task.

\subsection{Discussion to finalize with measurements}
\noindent\Pending. Insert the exact observed outcome for each research question here. Explain whether E2--E1 isolates a consistent morphology-specific gain, whether quality/diversity selection helped at matched sample counts, which classes improved or deteriorated, and whether changes persist across architectures and seeds. Discuss observed domain shift and computational costs, including unsuccessful runs and negative findings. Update only after the statistics, source audit and clinical annotation quality are checked.

\section{Ethics, Reproducibility and Conclusion}
\subsection{Data governance and model boundaries}
Only images with adequate legal authorization and an appropriate research-use basis should be retained. The image-source review must preserve original attribution and restrictions, and public redistribution is not implied merely because images appear in a Kaggle collection. Real and synthetic facial images can be identifiable or resemble real people. Patient identifiers, original photograph paths, sensitive image metadata, source images and private error galleries should therefore remain outside public GitHub. Synthetic figures must be explicitly labeled as synthetic, inspected for possible reproduction of training photographs, and checked against publication permissions.

The six-way classifier is a research model. Its outputs are broad appearance categories and should not be presented as confirmed diagnoses, risk scores or recommendations for treatment. Clinical verification by a qualified reviewer is necessary before asserting medical fidelity of generated lesions or of source labels. No ethical approval, institutional permission, clinical reader study or patient consent is claimed in this paper unless the appropriate documentation is actually obtained.

\subsection{Evidence ledger and reconstruction}
Each experiment is linked to a source-code commit, verified real split and label-audit fingerprint, model and adapter revision, selected image and prompt IDs, prompt-family partition, RNG seed, image transforms, optimizer schedule, training history, checkpoint checksum and environment details. The generation code writes per-image provenance and draft review CSVs; approved synthetic pools are combined with source-aware checks and then split without crossing prompt families. The classifiers retain checkpoint histories, six-way probabilities and real-validation selection records. Evaluation scripts compute the real-test metrics and paired contrasts; supplementary tables record variability across training seeds.

The full logs, checkpoint files and per-image review manifests can be stored in protected, versioned Kaggle outputs. GitHub contains the code, configuration templates, safe aggregate reports and references to verified run IDs. A human review is necessary before publishing an aggregate file so that original image paths or patient identifiers are not inadvertently disclosed. A CI pass verifies syntax and offline unit behavior; it does not establish that a GPU experiment ran successfully or that a medical claim is correct.

\subsection{Availability and assistance disclosure}
The project implementation and manuscript source are maintained in the DermSynReal repository at \url{https://github.com/shanmukhdatta/DermSynReal}. Versioned model and dataset assets may have separate permissions and download conditions. The authors will disclose the use of coding and writing assistance in the final submitted version in accordance with the venue's policy, distinguish automated drafting from scientific contribution, and verify all claims, references, tables and figures. Contributor roles and any acknowledgments require confirmation before submission.

\subsection{Conclusion}
DermSynReal provides a controlled way to examine whether synthetic facial photographs add useful information to a six-category real-image classifier. Its distinctive requirement is that the claim must survive diagnosis review, group-conscious data separation, generator-matched prompt comparisons and evaluation on the same untouched real photographs. The framework also makes a negative result informative: realistic images may fail to improve real-domain learning, and a gain from generic generation need not be a gain from clinical prompt detail.

\noindent\textbf{Final measured conclusion:} \Pending. Once the real-test comparisons have been executed and reviewed, replace this line with a concise numerical answer to the main questions and an appropriate uncertainty qualification. Do not change the prior description of the study to imply that experiments were run before they actually were.

\bibliographystyle{plainnat}
\bibliography{references}
\end{document}
